\section{Conclusion}
\label{sec:conclusion}

This study presents the successful development and implementation of an autonomous system for non-destructive inspection of rotating wind turbines, utilizing a UAV drone. Firstly, the article outlines the topic of research and provides an extensive examination of relevant literature, laying a solid theoretical foundation. Secondly, a presentation of the implemented model matching algorithm and the 3D model generation process is given, reflecting a progressive design approach to cater to the requirements of the inspection setting. The model matching algorithm works on a data-stream coming in from a LiDAR sensor mounted on board, resulting in a approximation of the relative wind turbine position with regards to the UAV. The 3D model generation process allows for a modular and fast re-purposing based on the turbine at hand. We then detailed the technical aspects of the parametrized trajectory generation algorithm and the camera triggering system to obtain high resolution images of the critical rotor blade parts. Finally, the proposed system underwent a series of tests using real-world data, validating its practicality and effectiveness.

The work can be further developed a number of future directions, such as the application of this methodology to other infrastructure, e.g. bridges and towers. This area of possible expansion of applicability underlines the versatility of the system. Additionally, a higher level of parametrization in the inspection process can be explored as another potential advancement, which could refine the system further and increase its efficiency. This research, therefore, contributes to the ongoing development of autonomous non-destructive inspections using UAVs.